% 2009 年福建卷理科第 6 题：程序框图。
% 来源 PDF：普通高考/2009/2009福建理.pdf，SHA-256 44d650687de50fff3efc1d37624b1272c9045e96d41491464f6713a708f202b9。
% 原页：第 1 页（A3 三栏）第 2 栏上部，书页码 515。
% 重绘依据：初值 S=2、n=1；循环体 S=1/(1-S)、n=2n；判断 S=2，是则输出 n，否则回到 S 赋值框。
\begin{tikzpicture}[semithick, every node/.style={inner sep=2pt, font=\small}, >=stealth]
  \node[draw,rounded corners] (start) at (0,0) {开始};
  \node[draw,rectangle] (s0) at (0,-1.1) {$S=2$};
  \node[draw,rectangle] (n0) at (0,-2.2) {$n=1$};
  \node[draw,rectangle] (s1) at (0,-3.3) {$S=\dfrac1{1-S}$};
  \node[draw,rectangle] (n1) at (0,-4.4) {$n=2n$};
  \node[draw,diamond,aspect=2] (test) at (0,-5.7) {$S=2$};
  \node[draw,rectangle] (out) at (0,-7.0) {输出 $n$};
  \node[draw,rounded corners] (end) at (0,-8.1) {结束};

  \draw[->] (start) -- (s0);
  \draw[->] (s0) -- (n0);
  \draw[->] (n0) -- (s1);
  \draw[->] (s1) -- (n1);
  \draw[->] (n1) -- (test);
  \draw[->] (test) -- node[right] {是} (out);
  \draw[->] (out) -- (end);
  \draw[->] (test.west) -- node[above] {否} (-1.9,-5.7) -- (-1.9,-3.3) -- (s1.west);
\end{tikzpicture}
